\begin{document}
\docTitle{Results guidelines}
Any result artifact produced from an experiment must be generated programmatically. 
This includes tables, plots, statistical tests, and any other results included in the experimental section of a document.
Each result artifact must be generated by a Python script that takes a single JSON file as input:
\begin{DocCode}
python3 script_name.py data.json
\end{DocCode}
Here, \code{script_name.py} and \code{data.json} are placeholders. Each artifact must use the actual script path and the actual input JSON path used to generate it.
Multiple artifacts may be generated from the same script and/or the same JSON input file.
The input JSON file must be created programmatically by a Python script named \code{collect_<json_name>.py}.
This collection script must gather the data from the relevant data sources and produce the JSON file used by the artifact-generation scripts.
It should not compute statics or aggregate results.
The final JSON should be a list of entries where
\begin{itemize}
    \item Each entry must include columns that uniquely identify each experiment.
    \item Each entry must include the location of the original data source.
    \item Each entry must include the commit hash associated with the data.
    \item Each entry must include the script used to obtain or process the data.
\end{itemize}
Each generated artifact must include provenance information. For every generated \code{.tex} table, a corresponding provenance file must be created:
\begin{DocCode}
<table_file_name>_provenance.json
\end{DocCode}
The provenance JSON must be a list with one entry for each table cell or result value. Each entry must include:
\begin{itemize}
\item the row identifier in the generated artifact;
\item the column identifier in the generated artifact;
\item the LaTeX value shown in the artifact;
\item the unique identifier of the source row used to produce that value;
\item the original source row data needed to reproduce or verify the value.
\end{itemize}
Generated tables must be stored under:

\begin{DocCode}
manual/generated/tables/<document_or_topic>/<table_name>.tex
\end{DocCode}

They must be imported into the document using the same path. For example:

\begin{DocCode}
\begin{table}[ht]
\centering
\caption{Caption text}
\label{tab:acc_main_experiment_table}
\input{manual/generated/tables/main/acc_main_experiment_table}
\end{table}
\end{DocCode}

Every generated table file must end with the following comments:

\begin{DocCode}
% This table was generated using python3 script_name.py data.json
% data.json was created using python3 collect_<json_name>.py
% Provenance of the data can be found at <table_file_name>_provenance.json
\end{DocCode}

The values in these comments must match the actual script names, input JSON file, collection script, and provenance file used to generate the table.

\subsection*{Pre result analysis}
Results can be analysed by automatically creating a \code{.tex} file

\end{document}
